\documentclass[12pt,a4paper]{article}
\usepackage[T2A]{fontenc}
\usepackage[utf8]{inputenc}
\usepackage[russian]{babel}
\usepackage{amsmath,amssymb,booktabs,array}
\usepackage[left=25mm,right=20mm,top=20mm,bottom=20mm]{geometry}
\usepackage{pgfplots}
\pgfplotsset{compat=1.18}
\usepackage[hidelinks]{hyperref}
\usepackage{float}
\setlength{\parindent}{1.25cm}
\setlength{\emergencystretch}{3em}
\linespread{1.12}
\newcommand{\ran}{\operatorname{ran}}
\newcommand{\rad}{\operatorname{rad}}
\newcommand{\midpoint}{\operatorname{mid}}
\newcommand{\dist}{\operatorname{dist}}
\newcommand{\hull}{\mathop{\square}}
\newcommand{\exprel}{\operatorname{exprel}}
\newcommand{\sinc}{\operatorname{sinc}}
\begin{document}
\begin{titlepage}
\setlength{\baselineskip}{15pt}
\begin{center}
Санкт-Петербургский политехнический университет Петра Великого\\
Физико-механический институт\\
Высшая школа прикладной математики и вычислительной физики
\vspace{30mm}

Направление подготовки\\
"01.03.02. Прикладная математика и информатика"
\vspace{27mm}

{\fontsize{20}{24}\selectfont Дисциплина "Интервальный анализ"}\\[5mm]
Отчет по лабораторной работе №2\\[2mm]
"По интервальному анализу"\\[4mm]
Вариант 23
\end{center}
\vspace{35mm}
\begin{flushright}
\begin{minipage}{38mm}
\raggedright
\textbf{Работу\\выполнил:}\\
Медведьков Р. А.\\
Группа:\\
5030102/30201
\end{minipage}
\end{flushright}
\vfill
\begin{center}Санкт-Петербург\end{center}
\end{titlepage}

\section*{Исходные данные}
Исследуются методы внешнего оценивания области значений: естественное
интервальное расширение, эквивалентные записи выражений, дифференциальная,
наклонная и бицентрированная формы. Состав работы соответствует пунктам
A--E задания от 17 сентября 2026 г.
Общая и индивидуальная функции одной переменной имеют вид
\begin{equation}\label{eq:input}
f_1(x)=x^3-3x^2+2,\quad X_1=[0,3];\qquad
f_2(x)=e^x-2x-1,\quad X_2=[0,2].
\end{equation}
Интервал $X_2$ выбран самостоятельно. Он содержит точку минимума
$x=\ln2$, а по разные стороны от неё функция меняет направление
монотонности. Для функции двух переменных в варианте 23 задано
\begin{equation}\label{eq:f3}
f_3(x,y)=x^2+y^2+\sin x+y\cos x,\qquad
X_3=X_{3,1}\times X_{3,2}=[-1,2]\times[-1,1].
\end{equation}
Все расчёты воспроизводятся файлом \texttt{lab2.py} без сторонних
библиотек. В отчёте использованы те же вычислительные выражения,
что и в программе. Полиномиальные результаты записаны точно,
если их десятичное представление конечно. Остальные значения
вычислены в двойной точности и округлены в таблицах для чтения.
После основных интервальных операций программа расширяет границы
с помощью \texttt{math.nextafter}. Трансцендентные функции вычисляются
стандартной библиотекой \texttt{math}, поэтому численные результаты
не являются сертификатом точного округления для любой платформы.
Математическое обоснование включений приведено в соответствующих
разделах; приближённые границы обозначены знаком $\approx$.

\section*{A. Точные области значений и графики}
Для общей функции $f'_1(x)=3x(x-2)$, поэтому достаточно проверить
концы отрезка и внутреннюю критическую точку $x=2$:
\[
f_1(0)=2,\qquad f_1(2)=-2,\qquad f_1(3)=2,
\qquad\boxed{\ran(f_1,X_1)=[-2,2]}.
\]
Для варианта 23 имеем
\[
f'_2(x)=e^x-2,\qquad f''_2(x)=e^x>0.
\]
Единственная критическая точка $\ln2\in[0,2]$ является строгим
минимумом. Значения на концах: $f_2(0)=0$, $f_2(2)=e^2-5>0$.
Следовательно,
\begin{equation}\label{eq:range2}
\boxed{\ran(f_2,X_2)=[1-2\ln2,e^2-5]
\approx[-0.3862943611,2.3890560989]}.
\end{equation}
У $f_2$ нет точек перегиба, однако наличие внутреннего экстремума
делает выбранный интервал содержательным для сравнения методов.

\begin{figure}[H]\centering
\begin{tikzpicture}
\begin{axis}[width=.92\linewidth,height=5.5cm,grid=both,
xlabel={$x$},ylabel={$f_1(x)$},xmin=0,xmax=3,domain=0:3,samples=121]
\addplot[blue,thick]{x^3-3*x^2+2};
\addplot[only marks,mark=*,red] coordinates {(0,2)(2,-2)(3,2)};
\end{axis}\end{tikzpicture}
\caption{Общая функция и точки достижения её границ на $[0,3]$.}
\end{figure}
\begin{figure}[H]\centering
\begin{tikzpicture}
\begin{axis}[width=.92\linewidth,height=5.5cm,grid=both,
xlabel={$x$},ylabel={$f_2(x)$},xmin=0,xmax=2,domain=0:2,samples=121]
\addplot[blue,thick]{exp(x)-2*x-1};
\addplot[only marks,mark=*,red] coordinates
{(0,0)(0.69314718056,-0.38629436112)(2,2.38905609893)};
\end{axis}\end{tikzpicture}
\caption{Функция варианта 23: минимум при $x=\ln2$ и максимум при $x=2$.}
\end{figure}

\section*{B. Интервальные оценки областей значений}
Интервал $F(X)$ называется внешней оценкой, если он содержит
$\ran(f,X)$. В пунктах B.1--B.2 степень $X^n$ вычисляется
повторным интервальным умножением.
В частности, $[-1,1]\cdot[-1,1]=[-1,1]$.
Отдельно будем обозначать через $Q(X)$ точное элементарное
расширение функции $x\mapsto x^2$; при $0\in X$ его нижняя граница
равна нулю. Эти операции различаются. Их роль в двумерном примере
будет указана явно, чтобы сравнение записей не скрывало смену операции.

\subsection*{B.1. Естественное расширение исходной записи}
Последовательная замена операций интервальными даёт
\[
F_{1,B1}=X_1^3-3X_1^2+2=[0,27]-[0,27]+2=[-25,29],
\]
\[
F_{2,B1}=e^{X_2}-2X_2-1=[1,e^2]-[0,4]-1=[-4,e^2-1].
\]
По отдельности границы слагаемых верны, но они могут соответствовать
разным аргументам. Их независимое комбинирование и создаёт избыточную
ширину оценки.

\subsection*{B.2. Естественное расширение эквивалентной записи}
Схема Горнера для общей функции даёт
\[
f_1(x)=x\bigl(x(x-3)\bigr)+2,
\quad X_1(X_1-3)=[-9,0],\quad F_{1,B2}=[-25,2].
\]
Во второй функции сгруппируем линейные члены:
\[
f_2(x)=e^x-2(x+1/2),\qquad
F_{2,B2}=[1,e^2]-2[1/2,5/2]=[-4,e^2-1].
\]
В данном случае преобразование не меняет результат: независимость
экспоненциального и линейного слагаемых сохраняется.
Это допустимое эквивалентное выражение, но оно не является улучшением.

\subsection*{B.3. Дифференциальная центрированная форма}
Пусть $D$ содержит все значения производной на $X$. Тогда при $c\in X$
\begin{equation}\label{eq:mv}
F_{\rm mv}(X,c)=f(c)+D(X-c)\supseteq\ran(f,X).
\end{equation}
Условие $c\in X$ обеспечивает применимость теоремы Лагранжа:
отрезок между центром и любой точкой $x\in X$ остаётся в области,
на которой оценена производная. Внешний центр также допустим,
если расширить эту область до $Y=\hull(X\cup\{c\})$ и использовать
$D\supseteq f'(Y)$. В нескольких переменных нужно охватить
все отрезки между центром и точками бруса, например расширенным брусом.
Оценивание производной только на исходном $X$ в таком случае
недостаточно. Для $f(t)=t^3$, $X=[0,1]$ и $c=2$ подстановка
$D=f'(X)=[0,3]$ даёт $8+[0,3][-2,-1]=[2,8]$,
не содержащий $f(0)=0$. При использовании $Y=[0,2]$
внешнее включение восстанавливается. Во всех последующих расчётах
центры выбраны внутри исследуемой области.

Для $f_1$ производная принимает минимум $-3$ при $x=1$ и максимум
$9$ при $x=3$. Для $f_2$ производная монотонно возрастает. Поэтому
используются точные диапазоны
\[
D_1=[-3,9],\qquad D_2=[-1,e^2-2]
\approx[-1,5.389056099].
\]
Центры выбираются в левом конце, середине и правом конце интервала.

\begin{center}
{\small
\begin{tabular}{ccc}
\toprule
Функция & Центр $c$ & $F_{\rm mv}(X,c)$\\\midrule
$f_1$ & $0$ & $[-7,\,29]$\\
$f_1$ & $1.5$ & $[-14.875,\,12.125]$\\
$f_1$ & $3$ & $[-25,\,11]$\\
$f_2$ & $0$ & $[-2,\,10.778112]$\\
$f_2$ & $1$ & $[-5.670774,\,5.107338]$\\
$f_2$ & $2$ & $[-8.389056,\,4.389056]$\\
\bottomrule
\end{tabular}}
\end{center}

Середина не гарантирует оптимальность среди всех возможных центров,
но здесь она даёт меньшую погрешность, чем оба конца, для обеих функций.

\subsection*{B.4. Наклонная центрированная форма}
Определим непрерывно продолженный наклон
\[
s(x,c)=\begin{cases}(f(x)-f(c))/(x-c),&x\ne c,\\f'(c),&x=c.\end{cases}
\]
Для интервала $S(X,c)\supseteq\{s(x,c):x\in X\}$ внешняя оценка равна
\begin{equation}\label{eq:sl}
F_{\rm sl}(X,c)=f(c)+S(X,c)(X-c).
\end{equation}
По теореме Лагранжа каждый точечный наклон при $c,x\in X$
совпадает с производной в промежуточной точке. Поэтому точный
диапазон наклонов содержится в точном диапазоне производной.
При согласованных внешних оценках $S(X,c)\subseteq D$
монотонность интервальных операций даёт
$F_{\rm sl}(X,c)\subseteq F_{\rm mv}(X,c)$.
Это объясняет возможное преимущество наклонной формы:
она учитывает секущие из выбранного центра, а не все касательные
на исходном интервале. В программе вложение проверяется для всех
шести рассматриваемых центров. Если же интервал наклонов вычислен
грубым интервальным выражением, такое вложение заранее не гарантируется.

Для полинома частное сокращается:
\[
s_1(x,c)=x^2+(c-3)x+c^2-3c.
\]
Его точный диапазон находится по концам $X_1$ и вершине параболы
$x=(3-c)/2$, если она лежит в $X_1$. Поэтому дополнительных
интервальных переоценок при вычислении $S_1$ здесь нет.

Для экспоненциальной функции получаем
\begin{equation}\label{eq:exprel}
s_2(x,c)=e^c\exprel(x-c)-2,\qquad
\exprel(t)=\begin{cases}(e^t-1)/t,&t\ne0,\\1,&t=0.\end{cases}
\end{equation}
Представление $\exprel(t)=\int_0^1e^{ut}\,du$ показывает, что
$\exprel$ строго возрастает на всей вещественной оси. Поэтому точные
границы наклонов задаются концами $X_2-c$.
Программа использует \texttt{math.expm1}, чтобы при малом $t$ избежать
потери точности при вычитании близких чисел.
Нуль в $X_2-c$ не является разрывом продолженного наклона.
В точке $t=0$ применяется предельное значение $\exprel(0)=1$,
а на всём интервале --- монотонность продолженной функции.
Обычное интервальное деление $(e^{X_2-c}-1)/(X_2-c)$ здесь
недопустимо, поскольку знаменатель содержит нуль.

Аналогичный приём используется в наклонах тригонометрических
слагаемых. Например,
$\sin(2x)-\sin(2c)=2\cos(x+c)\sin(x-c)$ приводит к функции
$\sinc(t)=\sin t/t$, доопределяемой значением $\sinc(0)=1$.
Нули числителя при $t=k\pi\ne0$ особенностей не создают:
там знаменатель ненулевой и отношение равно нулю.
При оценивании $\sinc$ на широком интервале необходимо учитывать
её внутренние экстремумы вместе с концами и точкой непрерывного
продолжения. В варианте 23 для одномерной функции используется
именно экспоненциальное продолжение~\eqref{eq:exprel};
синус в $f_3$ оценивается непосредственно с учётом его экстремумов.

\begin{center}
{\small
\begin{tabular}{cccc}
\toprule
Функция & $c$ & $S(X,c)$ & $F_{\rm sl}(X,c)$\\\midrule
$f_1$ & $0$ & $[-2.25,\,0]$ & $[-4.75,\,2]$\\
$f_1$ & $1.5$ & $[-2.8125,\,2.25]$ & $[-5.59375,\,2.84375]$\\
$f_1$ & $3$ & $[0,\,9]$ & $[-25,\,2]$\\
$f_2$ & $0$ & $[-1,\,1.194528]$ & $[-2,\,2.389056]$\\
$f_2$ & $1$ & $[-0.281718,\,2.670774]$ & $[-2.952492,\,2.389056]$\\
$f_2$ & $2$ & $[1.194528,\,5.389056]$ & $[-8.389056,\,2.389056]$\\
\bottomrule
\end{tabular}}
\end{center}

Например, при $c=0$ для первой функции
$S_1=[-9/4,0]$, и $F_{\rm sl}=2+[-9/4,0][0,3]=[-19/4,2]$.
Для второй функции тот же левый центр сохраняет точную верхнюю
границу $e^2-5$, хотя нижняя граница $-2$ остаётся избыточной.

\subsection*{B.5. Бицентрированная среднезначная форма}
При фиксированном интервале производной используем центры Бауманна:
\begin{equation}\label{eq:baumann}
p=\max\{-1,\min\{1,\midpoint D/\rad D\}\},\quad
c^-=\midpoint X-p\rad X,\quad c^+=\midpoint X+p\rad X.
\end{equation}
Если $\rad D=0$, можно взять $p=0$. Бицентрированная оценка получается
пересечением двух внешних интервалов:
\[
F_{\rm bic}=F_{\rm mv}(X,c^-)\cap F_{\rm mv}(X,c^+).
\]
Для $f_1$: $p=1/2$, $c^-=3/4$, $c^+=9/4$,
\[
F_{\rm mv}(X_1,3/4)=[-6.015625,20.984375],\quad
F_{\rm mv}(X_1,9/4)=[-22.046875,4.953125],
\]
\[
F_{1,B5}=[-6.015625,4.953125].
\]
Для $f_2$ имеем
\[
p=\frac{e^2-3}{e^2-1},\qquad c^-=1-p\approx0.313035286,
\qquad c^+=1+p\approx1.686964715.
\]
Численные среднезначные оценки в этих точках:
\[
F_{\rm mv}(X_2,c^-)\approx[-1.945465500,8.832646698],
\]
\[
F_{\rm mv}(X_2,c^+)\approx[-8.062020939,2.716091259],
\]
\[
F_{2,B5}\approx[-1.945465500,2.716091259].
\]
Оба интервала содержат область значений, поэтому её содержит
и их непустое пересечение.
Правило Бауманна относится к среднезначной форме с фиксированным
интервалом производной. Оно не гарантирует оптимальность этих же
центров для наклонной формы, где $S(X,c)$ зависит от самого $c$.
Это различие видно и по результатам $f_1$: наклонная оценка
при левом центре имеет расстояние $2.75$, тогда как пересечение
форм Бауманна даёт $4.015625$.

\section*{C. Оценка точности}
\subsection*{Расстояние по Хаусдорфу для всех оценок}
Для двух замкнутых интервалов справедливо
\begin{equation}\label{eq:hausdorff}
\dist([a,b],[m,M])=\max\{|a-m|,|b-M|\}.
\end{equation}
Если первый интервал внешний, модули можно заменить на $m-a$ и $b-M$.
Все расстояния ниже вычислены по неокруглённым значениям программы.
Числа для неполиномиальных функций приведены приближённо.

\begin{center}
{\small
\begin{tabular}{lcc}
\toprule
Метод и центр & $\dist(F_1,\ran f_1)$ & $\dist(F_2,\ran f_2)$\\\midrule
B.1, исходная запись & $27$ & $4$\\
B.2, преобразованная & $23$ & $4$\\
B.3, левый центр & $27$ & $8.389056$\\
B.3, середина & $12.875$ & $5.28448$\\
B.3, правый центр & $23$ & $8.002762$\\
B.4, левый центр & $2.75$ & $1.613706$\\
B.4, середина & $3.59375$ & $2.566198$\\
B.4, правый центр & $23$ & $8.002762$\\
B.5, два центра & $4.015625$ & $1.559171$\\
\bottomrule
\end{tabular}}
\end{center}


\subsection*{C.1. Константы Липшица по дереву Канторовича}
Сначала найдём константы самих точечных функций:
\[
L_{f_1}=\sup_{[0,3]}|f'_1|=9,\qquad
L_{f_2}=\sup_{[0,2]}|e^x-2|=e^2-2\approx5.389056099.
\]
Они гарантируют $\rad\ran(f,X)\leq L_f\rad X$.
Однако естественное расширение способно оказаться шире точного
диапазона. Для оценки его полуширины используем константу
вычисляемого выражения $L_{\rm expr}$, полученную по дереву Канторовича.
Для суммы и произведения применяются правила
\[
L_{g\pm h}\leq L_g+L_h,\qquad
L_{gh}\leq |G(X)|L_h+|H(X)|L_g,
\quad |[a,b]|=\max(|a|,|b|).
\]
У переменной константа равна единице, у числового узла --- нулю.
Оценивание узлов на всём исходном интервале даёт допустимые, хотя
часто завышенные, константы для соответствующей интервальной записи.

Для $f_1$ получаем $L_{x^2}=6$, $L_{x^3}=27$,
поэтому для исходной записи $L_{1,B1}=27+3\cdot6=45$.
Для Горнера последовательно имеем
\[
L_{x-3}=1,\quad L_{x(x-3)}=3+3=6,\quad
L_{x(x(x-3))}=9+3\cdot6=27.
\]
Для обеих записей $f_2$:
\[
L_{e^x}=e^2,\qquad L_{2x}=L_{2(x+1/2)}=2,
\quad L_{2,B1}=L_{2,B2}=e^2+2\approx9.389056099.
\]
Совпадение чисел в последней строке объясняется тем, что сдвиг
линейного аргумента не меняет его константу.

\subsection*{C.2--C.3. Теоретические границы и фактические полуширины}
Сравниваем фактическую полуширину естественной оценки с границей
\begin{equation}\label{eq:lip}
\rad F(X)\leq L_{\rm expr}\rad X,\qquad
\rad[a,b]=(b-a)/2.
\end{equation}

\begin{center}
{\small
\begin{tabular}{lccc}
\toprule
Запись & $L_{\rm expr}$ & $\rad F$ & $L_{\rm expr}\rad X$\\\midrule
$f_1$, B.1 & $45$ & $27$ & $67.5$\\
$f_1$, B.2 & $27$ & $13.5$ & $40.5$\\
$f_2$, B.1 & $9.389056$ & $5.194528$ & $9.389056$\\
$f_2$, B.2 & $9.389056$ & $5.194528$ & $9.389056$\\
\bottomrule
\end{tabular}}
\end{center}

Для $f_1$ подстановка только $L_{f_1}=9$ в оценку естественной
формы неверна: $27>9\cdot1.5=13.5$.
Противоречия с липшицевостью самой функции нет, поскольку её точный
диапазон имеет радиус $2$. Для $f_2$ неравенство с $L_{f_2}$
в данном случае также выполняется: $5.194528050\leq5.389056099$,
но это частный результат, а не общее основание для оценки
произвольного естественного расширения.
Расстояние~\eqref{eq:hausdorff} характеризует избыточность внешней
оценки, тогда как~\eqref{eq:lip} ограничивает всю её полуширину.
Эти величины нельзя отождествлять.

\section*{D. Сравнение методов и выводы}
Вычисления показывают, что точность зависит одновременно от записи
функции и способа учёта зависимости повторных вхождений аргумента.
Для $f_1$ переход к Горнеру сохраняет избыточную нижнюю границу,
но делает верхнюю точной: расстояние уменьшается с $27$ до $23$.
У $f_2$ перегруппировка линейных членов не меняет ни одной границы.

Среди выбранных центров дифференциальной формы середина даёт лучшие
результаты для обеих функций. В наклонной форме лучшим среди
рассмотренных оказывается левый центр: расстояния равны $2.75$
для $f_1$ и примерно $1.613706$ для $f_2$.
Это различие связано с тем, что диапазон секущих зависит от центра.

Два центра Бауманна сокращают переоценку по сравнению с проверенными
одиночными среднезначными формами. Для $f_1$ их пересечение уступает
лучшему наклонному результату. Для $f_2$, напротив, бицентрированная
оценка имеет наименьшее расстояние среди всех вычисленных:
примерно $1.559171$ против $1.613706$ у лучшей наклонной.
Следовательно, ни одна из этих форм не оказалась лучшей сразу для
обеих функций. Липшицевские границы по дереву выражения выполнены,
но заметно превышают фактические полуширины.

\section*{E. Функция двух переменных}
\subsection*{E.1. График поверхности}
На рисунках приведены поверхность и линии уровня на всём заданном
брусе. Углы тригонометрических функций в математических формулах
заданы в радианах; при построении средствами PGFPlots выполняется
необходимое преобразование единиц.
\begin{figure}[H]\centering
\begin{tikzpicture}
\begin{axis}[width=.92\linewidth,height=7cm,view={125}{28},
xlabel={$x$},ylabel={$y$},zlabel={$f_3(x,y)$},
domain=-1:2,y domain=-1:1,samples=27,samples y=23,colormap/viridis,
colorbar,point meta min=-.45,point meta max=6.4]
\addplot3[surf,shader=interp] {x^2+y^2+sin(deg(x))+y*cos(deg(x))};
\end{axis}\end{tikzpicture}
\caption{Поверхность $z=x^2+y^2+\sin x+y\cos x$.}
\end{figure}
\begin{figure}[H]\centering
\begin{tikzpicture}
\begin{axis}[width=.88\linewidth,height=7cm,xmin=-1,xmax=2,ymin=-1,ymax=1,
xlabel={$x$},ylabel={$y$},grid=major,unbounded coords=jump,
legend style={font=\scriptsize,at={(1.02,1)},anchor=north west}]
\addplot[no marks,thin,violet] coordinates {
(-0.64000,-0.70925) (-0.66197,-0.66860) (nan,nan)
(-0.66197,-0.66860) (-0.67603,-0.64706) (nan,nan)
(-0.67603,-0.64706) (-0.69775,-0.58603) (nan,nan)
(-0.69775,-0.58603) (-0.72393,-0.52941) (nan,nan)
(-0.72393,-0.52941) (-0.72835,-0.49838) (nan,nan)
(-0.72835,-0.49838) (-0.74916,-0.41176) (nan,nan)
(-0.74916,-0.41176) (-0.74789,-0.39989) (nan,nan)
(-0.74789,-0.39989) (-0.74546,-0.29412) (nan,nan)
(-0.74546,-0.29412) (-0.73601,-0.27059) (nan,nan)
(-0.73601,-0.27059) (-0.70407,-0.17647) (nan,nan)
(-0.70407,-0.17647) (-0.64000,-0.09179) (nan,nan)
(-0.52000,-0.83571) (-0.55521,-0.79923) (nan,nan)
(-0.55521,-0.79923) (-0.59448,-0.76471) (nan,nan)
(-0.59448,-0.76471) (-0.61350,-0.73873) (nan,nan)
(-0.61350,-0.73873) (-0.64000,-0.70925) (nan,nan)
(-0.58756,-0.05882) (-0.64000,-0.09179) (nan,nan)
(-0.58756,-0.05882) (-0.52000,-0.03109) (nan,nan)
(-0.40000,-0.89647) (-0.40962,-0.89179) (nan,nan)
(-0.40962,-0.89179) (-0.43263,-0.88235) (nan,nan)
(-0.43263,-0.88235) (-0.48777,-0.85076) (nan,nan)
(-0.48777,-0.85076) (-0.52000,-0.83571) (nan,nan)
(-0.40000,-0.02610) (-0.49176,-0.03114) (nan,nan)
(-0.49176,-0.03114) (-0.52000,-0.03109) (nan,nan)
(-0.28000,-0.90124) (-0.29786,-0.89986) (nan,nan)
(-0.29786,-0.89986) (-0.40000,-0.89647) (nan,nan)
(-0.28000,-0.05775) (-0.37456,-0.03388) (nan,nan)
(-0.37456,-0.03388) (-0.40000,-0.02610) (nan,nan)
(-0.22926,-0.88235) (-0.28000,-0.90124) (nan,nan)
(-0.22926,-0.88235) (-0.16000,-0.84637) (nan,nan)
(-0.16000,-0.14084) (-0.20997,-0.10782) (nan,nan)
(-0.20997,-0.10782) (-0.27803,-0.05882) (nan,nan)
(-0.27803,-0.05882) (-0.27931,-0.05815) (nan,nan)
(-0.27931,-0.05815) (-0.28000,-0.05775) (nan,nan)
(-0.08497,-0.76471) (-0.16000,-0.84637) (nan,nan)
(-0.08497,-0.76471) (-0.04000,-0.68644) (nan,nan)
(-0.04000,-0.30855) (-0.04407,-0.29811) (nan,nan)
(-0.04407,-0.29811) (-0.04553,-0.29412) (nan,nan)
(-0.04553,-0.29412) (-0.09176,-0.22721) (nan,nan)
(-0.09176,-0.22721) (-0.12551,-0.17647) (nan,nan)
(-0.12551,-0.17647) (-0.14244,-0.15925) (nan,nan)
(-0.14244,-0.15925) (-0.16000,-0.14084) (nan,nan)
(-0.02456,-0.64706) (-0.04000,-0.68644) (nan,nan)
(-0.02456,-0.64706) (-0.02145,-0.62888) (nan,nan)
(-0.02145,-0.62888) (-0.00466,-0.52941) (nan,nan)
(-0.00466,-0.52941) (-0.00640,-0.49647) (nan,nan)
(-0.00640,-0.49647) (-0.01110,-0.41176) (nan,nan)
(-0.01110,-0.41176) (-0.01728,-0.38949) (nan,nan)
(-0.01728,-0.38949) (-0.04000,-0.30855) (nan,nan)
};
\addlegendentry{$z=-0.25$}
\addplot[no marks,thin,blue] coordinates {
(-0.88000,-0.62838) (-0.90280,-0.55176) (nan,nan)
(-0.90280,-0.55176) (-0.91107,-0.52941) (nan,nan)
(-0.91107,-0.52941) (-0.92422,-0.45512) (nan,nan)
(-0.92422,-0.45512) (-0.93493,-0.41176) (nan,nan)
(-0.93493,-0.41176) (-0.93752,-0.35051) (nan,nan)
(-0.93752,-0.35051) (-0.94364,-0.29412) (nan,nan)
(-0.94364,-0.29412) (-0.93689,-0.23225) (nan,nan)
(-0.93689,-0.23225) (-0.93418,-0.17647) (nan,nan)
(-0.93418,-0.17647) (-0.90922,-0.08747) (nan,nan)
(-0.90922,-0.08747) (-0.90269,-0.05882) (nan,nan)
(-0.90269,-0.05882) (-0.88000,-0.01118) (nan,nan)
(-0.76000,-0.85236) (-0.79224,-0.79632) (nan,nan)
(-0.79224,-0.79632) (-0.81329,-0.76471) (nan,nan)
(-0.81329,-0.76471) (-0.83380,-0.71941) (nan,nan)
(-0.83380,-0.71941) (-0.87296,-0.64706) (nan,nan)
(-0.87296,-0.64706) (-0.87462,-0.64179) (nan,nan)
(-0.87462,-0.64179) (-0.88000,-0.62838) (nan,nan)
(-0.83127,0.05882) (-0.88000,-0.01118) (nan,nan)
(-0.83127,0.05882) (-0.76000,0.12676) (nan,nan)
(-0.64000,-0.99046) (-0.68994,-0.93131) (nan,nan)
(-0.68994,-0.93131) (-0.73702,-0.88235) (nan,nan)
(-0.73702,-0.88235) (-0.74639,-0.86901) (nan,nan)
(-0.74639,-0.86901) (-0.76000,-0.85236) (nan,nan)
(-0.66856,0.17647) (-0.76000,0.12676) (nan,nan)
(-0.66856,0.17647) (-0.64000,0.18818) (nan,nan)
(-0.62817,-1.00000) (-0.63448,-0.99458) (nan,nan)
(-0.63448,-0.99458) (-0.64000,-0.99046) (nan,nan)
(-0.52000,0.20800) (-0.62626,0.18994) (nan,nan)
(-0.62626,0.18994) (-0.64000,0.18818) (nan,nan)
(-0.40000,0.20215) (-0.49025,0.20564) (nan,nan)
(-0.49025,0.20564) (-0.52000,0.20800) (nan,nan)
(-0.28000,0.17414) (-0.28191,0.17459) (nan,nan)
(-0.28191,0.17459) (-0.28867,0.17647) (nan,nan)
(-0.28867,0.17647) (-0.37928,0.19679) (nan,nan)
(-0.37928,0.19679) (-0.40000,0.20215) (nan,nan)
(-0.16000,0.11787) (-0.20101,0.13626) (nan,nan)
(-0.20101,0.13626) (-0.28000,0.17414) (nan,nan)
(-0.04000,0.03496) (-0.05349,0.04560) (nan,nan)
(-0.05349,0.04560) (-0.06980,0.05882) (nan,nan)
(-0.06980,0.05882) (-0.12410,0.09402) (nan,nan)
(-0.12410,0.09402) (-0.16000,0.11787) (nan,nan)
(-0.00451,-1.00000) (0.08000,-0.89892) (nan,nan)
(0.08000,-0.09972) (0.06260,-0.07589) (nan,nan)
(0.06260,-0.07589) (0.05007,-0.05882) (nan,nan)
(0.05007,-0.05882) (0.00644,-0.01329) (nan,nan)
(0.00644,-0.01329) (-0.04000,0.03496) (nan,nan)
(0.09053,-0.88235) (0.08000,-0.89892) (nan,nan)
(0.09053,-0.88235) (0.09986,-0.86289) (nan,nan)
(0.09986,-0.86289) (0.14624,-0.76471) (nan,nan)
(0.14624,-0.76471) (0.17566,-0.67092) (nan,nan)
(0.17566,-0.67092) (0.18294,-0.64706) (nan,nan)
(0.18294,-0.64706) (0.19993,-0.52949) (nan,nan)
(0.19993,-0.52949) (0.19993,-0.52941) (nan,nan)
(0.19993,-0.52941) (0.19663,-0.41507) (nan,nan)
(0.19663,-0.41507) (0.19649,-0.41176) (nan,nan)
(0.19649,-0.41176) (0.17687,-0.31680) (nan,nan)
(0.17687,-0.31680) (0.17183,-0.29412) (nan,nan)
(0.17183,-0.29412) (0.14649,-0.22893) (nan,nan)
(0.14649,-0.22893) (0.12516,-0.17647) (nan,nan)
(0.12516,-0.17647) (0.10887,-0.14816) (nan,nan)
(0.10887,-0.14816) (0.08000,-0.09972) (nan,nan)
};
\addlegendentry{$z=0$}
\addplot[no marks,thin,cyan!70!black] coordinates {
(-0.94364,-1.00000) (-0.96437,-0.96507) (nan,nan)
(-0.96437,-0.96507) (-1.00000,-0.91191) (nan,nan)
(-0.94724,0.41176) (-1.00000,0.37120) (nan,nan)
(-0.94724,0.41176) (-0.88000,0.45259) (nan,nan)
(-0.76000,0.49778) (-0.81399,0.47648) (nan,nan)
(-0.81399,0.47648) (-0.88000,0.45259) (nan,nan)
(-0.64000,0.51948) (-0.65294,0.51673) (nan,nan)
(-0.65294,0.51673) (-0.76000,0.49778) (nan,nan)
(-0.52000,0.52210) (-0.52790,0.52167) (nan,nan)
(-0.52790,0.52167) (-0.64000,0.51948) (nan,nan)
(-0.40000,0.50875) (-0.41940,0.51040) (nan,nan)
(-0.41940,0.51040) (-0.52000,0.52210) (nan,nan)
(-0.28000,0.48151) (-0.32024,0.48996) (nan,nan)
(-0.32024,0.48996) (-0.40000,0.50875) (nan,nan)
(-0.16000,0.44167) (-0.22726,0.46347) (nan,nan)
(-0.22726,0.46347) (-0.28000,0.48151) (nan,nan)
(-0.04000,0.38678) (-0.05707,0.39503) (nan,nan)
(-0.05707,0.39503) (-0.09098,0.41176) (nan,nan)
(-0.09098,0.41176) (-0.13894,0.43241) (nan,nan)
(-0.13894,0.43241) (-0.16000,0.44167) (nan,nan)
(0.08000,0.31409) (0.01846,0.35143) (nan,nan)
(0.01846,0.35143) (-0.04000,0.38678) (nan,nan)
(0.20000,0.21593) (0.15747,0.25242) (nan,nan)
(0.15747,0.25242) (0.10768,0.29412) (nan,nan)
(0.10768,0.29412) (0.09180,0.30569) (nan,nan)
(0.09180,0.30569) (0.08000,0.31409) (nan,nan)
(0.32000,0.07886) (0.27607,0.13340) (nan,nan)
(0.27607,0.13340) (0.23973,0.17647) (nan,nan)
(0.23973,0.17647) (0.22021,0.19628) (nan,nan)
(0.22021,0.19628) (0.20000,0.21593) (nan,nan)
(0.33567,-1.00000) (0.34933,-0.97125) (nan,nan)
(0.34933,-0.97125) (0.39049,-0.88235) (nan,nan)
(0.39049,-0.88235) (0.43033,-0.77418) (nan,nan)
(0.43033,-0.77418) (0.43369,-0.76471) (nan,nan)
(0.43369,-0.76471) (0.44000,-0.74016) (nan,nan)
(0.44000,-0.16306) (0.41104,-0.08721) (nan,nan)
(0.41104,-0.08721) (0.39919,-0.05882) (nan,nan)
(0.39919,-0.05882) (0.37214,-0.00770) (nan,nan)
(0.37214,-0.00770) (0.33424,0.05882) (nan,nan)
(0.33424,0.05882) (0.32853,0.06718) (nan,nan)
(0.32853,0.06718) (0.32000,0.07886) (nan,nan)
(0.46158,-0.64706) (0.44000,-0.74016) (nan,nan)
(0.46158,-0.64706) (0.46532,-0.62224) (nan,nan)
(0.46532,-0.62224) (0.47716,-0.52941) (nan,nan)
(0.47716,-0.52941) (0.47883,-0.49134) (nan,nan)
(0.47883,-0.49134) (0.48023,-0.41176) (nan,nan)
(0.48023,-0.41176) (0.47771,-0.37479) (nan,nan)
(0.47771,-0.37479) (0.46976,-0.29412) (nan,nan)
(0.46976,-0.29412) (0.46516,-0.26945) (nan,nan)
(0.46516,-0.26945) (0.44458,-0.17647) (nan,nan)
(0.44458,-0.17647) (0.44351,-0.17303) (nan,nan)
(0.44351,-0.17303) (0.44000,-0.16306) (nan,nan)
};
\addlegendentry{$z=0.5$}
\addplot[no marks,thin,teal] coordinates {
(-0.88000,0.72777) (-0.94100,0.70490) (nan,nan)
(-0.94100,0.70490) (-1.00000,0.68455) (nan,nan)
(-0.76000,0.75172) (-0.77724,0.74781) (nan,nan)
(-0.77724,0.74781) (-0.88000,0.72777) (nan,nan)
(-0.64000,0.75994) (-0.64540,0.75941) (nan,nan)
(-0.64540,0.75941) (-0.76000,0.75172) (nan,nan)
(-0.52000,0.75509) (-0.52968,0.75522) (nan,nan)
(-0.52968,0.75522) (-0.64000,0.75994) (nan,nan)
(-0.40000,0.73912) (-0.42341,0.74175) (nan,nan)
(-0.42341,0.74175) (-0.52000,0.75509) (nan,nan)
(-0.28000,0.71340) (-0.32342,0.72214) (nan,nan)
(-0.32342,0.72214) (-0.40000,0.73912) (nan,nan)
(-0.16000,0.67884) (-0.22802,0.69802) (nan,nan)
(-0.22802,0.69802) (-0.28000,0.71340) (nan,nan)
(-0.04000,0.63472) (-0.04900,0.63823) (nan,nan)
(-0.04900,0.63823) (-0.07124,0.64706) (nan,nan)
(-0.07124,0.64706) (-0.13634,0.67025) (nan,nan)
(-0.13634,0.67025) (-0.16000,0.67884) (nan,nan)
(0.08000,0.57804) (0.03252,0.60051) (nan,nan)
(0.03252,0.60051) (-0.04000,0.63472) (nan,nan)
(0.20000,0.50973) (0.18779,0.51744) (nan,nan)
(0.18779,0.51744) (0.16838,0.52941) (nan,nan)
(0.16838,0.52941) (0.11193,0.56072) (nan,nan)
(0.11193,0.56072) (0.08000,0.57804) (nan,nan)
(0.32000,0.42373) (0.25821,0.46884) (nan,nan)
(0.25821,0.46884) (0.20000,0.50973) (nan,nan)
(0.44000,0.31141) (0.38882,0.36159) (nan,nan)
(0.38882,0.36159) (0.33473,0.41176) (nan,nan)
(0.33473,0.41176) (0.32676,0.41839) (nan,nan)
(0.32676,0.41839) (0.32000,0.42373) (nan,nan)
(0.56000,0.15297) (0.55126,0.16790) (nan,nan)
(0.55126,0.16790) (0.54581,0.17647) (nan,nan)
(0.54581,0.17647) (0.50149,0.23675) (nan,nan)
(0.50149,0.23675) (0.45604,0.29412) (nan,nan)
(0.45604,0.29412) (0.44855,0.30250) (nan,nan)
(0.44855,0.30250) (0.44000,0.31141) (nan,nan)
(0.56094,-1.00000) (0.56159,-0.99844) (nan,nan)
(0.56159,-0.99844) (0.60842,-0.88235) (nan,nan)
(0.60842,-0.88235) (0.63252,-0.81126) (nan,nan)
(0.63252,-0.81126) (0.64743,-0.76471) (nan,nan)
(0.64743,-0.76471) (0.67641,-0.65058) (nan,nan)
(0.67641,-0.65058) (0.67722,-0.64706) (nan,nan)
(0.67722,-0.64706) (0.68000,-0.63016) (nan,nan)
(0.68000,-0.14902) (0.66180,-0.07667) (nan,nan)
(0.66180,-0.07667) (0.65657,-0.05882) (nan,nan)
(0.65657,-0.05882) (0.63074,0.01053) (nan,nan)
(0.63074,0.01053) (0.61047,0.05882) (nan,nan)
(0.61047,0.05882) (0.59383,0.09199) (nan,nan)
(0.59383,0.09199) (0.56000,0.15297) (nan,nan)
(0.69561,-0.52941) (0.68000,-0.63016) (nan,nan)
(0.69561,-0.52941) (0.69721,-0.51254) (nan,nan)
(0.69721,-0.51254) (0.70382,-0.41176) (nan,nan)
(0.70382,-0.41176) (0.70390,-0.38833) (nan,nan)
(0.70390,-0.38833) (0.70111,-0.29412) (nan,nan)
(0.70111,-0.29412) (0.69933,-0.27516) (nan,nan)
(0.69933,-0.27516) (0.68637,-0.17647) (nan,nan)
(0.68637,-0.17647) (0.68534,-0.17124) (nan,nan)
(0.68534,-0.17124) (0.68000,-0.14902) (nan,nan)
};
\addlegendentry{$z=1$}
\addplot[no marks,thin,green!50!black] coordinates {
(0.08000,0.97113) (0.05838,0.97880) (nan,nan)
(0.05838,0.97880) (-0.00155,1.00000) (nan,nan)
(0.20000,0.92358) (0.14466,0.94575) (nan,nan)
(0.14466,0.94575) (0.08000,0.97113) (nan,nan)
(0.32000,0.86961) (0.31138,0.87390) (nan,nan)
(0.31138,0.87390) (0.29373,0.88235) (nan,nan)
(0.29373,0.88235) (0.22925,0.91103) (nan,nan)
(0.22925,0.91103) (0.20000,0.92358) (nan,nan)
(0.44000,0.80539) (0.38972,0.83306) (nan,nan)
(0.38972,0.83306) (0.32000,0.86961) (nan,nan)
(0.56000,0.73042) (0.53952,0.74462) (nan,nan)
(0.53952,0.74462) (0.50872,0.76471) (nan,nan)
(0.50872,0.76471) (0.46624,0.79043) (nan,nan)
(0.46624,0.79043) (0.44000,0.80539) (nan,nan)
(0.68000,0.64040) (0.67645,0.64357) (nan,nan)
(0.67645,0.64357) (0.67224,0.64706) (nan,nan)
(0.67224,0.64706) (0.60924,0.69534) (nan,nan)
(0.60924,0.69534) (0.56000,0.73042) (nan,nan)
(0.80000,0.52546) (0.79817,0.52762) (nan,nan)
(0.79817,0.52762) (0.79652,0.52941) (nan,nan)
(0.79652,0.52941) (0.73865,0.58691) (nan,nan)
(0.73865,0.58691) (0.68000,0.64040) (nan,nan)
(0.90844,-1.00000) (0.92000,-0.97166) (nan,nan)
(0.92000,0.36525) (0.90224,0.39435) (nan,nan)
(0.90224,0.39435) (0.89039,0.41176) (nan,nan)
(0.89039,0.41176) (0.85206,0.46281) (nan,nan)
(0.85206,0.46281) (0.80000,0.52546) (nan,nan)
(0.95509,-0.88235) (0.92000,-0.97166) (nan,nan)
(0.95509,-0.88235) (0.97351,-0.82990) (nan,nan)
(0.97351,-0.82990) (0.99505,-0.76471) (nan,nan)
(0.99505,-0.76471) (1.02448,-0.66228) (nan,nan)
(1.02448,-0.66228) (1.02849,-0.64706) (nan,nan)
(1.02849,-0.64706) (1.04000,-0.59464) (nan,nan)
(1.04000,0.08951) (1.01982,0.15668) (nan,nan)
(1.01982,0.15668) (1.01280,0.17647) (nan,nan)
(1.01280,0.17647) (0.98609,0.24127) (nan,nan)
(0.98609,0.24127) (0.96106,0.29412) (nan,nan)
(0.96106,0.29412) (0.94696,0.32055) (nan,nan)
(0.94696,0.32055) (0.92000,0.36525) (nan,nan)
(1.05399,-0.52941) (1.04000,-0.59464) (nan,nan)
(1.05399,-0.52941) (1.05686,-0.51289) (nan,nan)
(1.05686,-0.51289) (1.07142,-0.41176) (nan,nan)
(1.07142,-0.41176) (1.07489,-0.37756) (nan,nan)
(1.07489,-0.37756) (1.08061,-0.29412) (nan,nan)
(1.08061,-0.29412) (1.08160,-0.25333) (nan,nan)
(1.08160,-0.25333) (1.08063,-0.17647) (nan,nan)
(1.08063,-0.17647) (1.07843,-0.13879) (nan,nan)
(1.07843,-0.13879) (1.07042,-0.05882) (nan,nan)
(1.07042,-0.05882) (1.06658,-0.03277) (nan,nan)
(1.06658,-0.03277) (1.04875,0.05882) (nan,nan)
(1.04875,0.05882) (1.04706,0.06575) (nan,nan)
(1.04706,0.06575) (1.04000,0.08951) (nan,nan)
};
\addlegendentry{$z=2$}
\addplot[no marks,thin,olive] coordinates {
(0.80000,0.97904) (0.78698,0.98723) (nan,nan)
(0.78698,0.98723) (0.76499,1.00000) (nan,nan)
(0.92000,0.90040) (0.86019,0.94136) (nan,nan)
(0.86019,0.94136) (0.80000,0.97904) (nan,nan)
(1.04000,0.80377) (0.99762,0.84081) (nan,nan)
(0.99762,0.84081) (0.94517,0.88235) (nan,nan)
(0.94517,0.88235) (0.93088,0.89302) (nan,nan)
(0.93088,0.89302) (0.92000,0.90040) (nan,nan)
(1.16000,0.68047) (1.12040,0.72588) (nan,nan)
(1.12040,0.72588) (1.08265,0.76471) (nan,nan)
(1.08265,0.76471) (1.06114,0.78543) (nan,nan)
(1.06114,0.78543) (1.04000,0.80377) (nan,nan)
(1.19679,-1.00000) (1.22347,-0.93778) (nan,nan)
(1.22347,-0.93778) (1.24623,-0.88235) (nan,nan)
(1.24623,-0.88235) (1.28000,-0.79269) (nan,nan)
(1.28000,0.50622) (1.27209,0.52166) (nan,nan)
(1.27209,0.52166) (1.26754,0.52941) (nan,nan)
(1.26754,0.52941) (1.22608,0.59420) (nan,nan)
(1.22608,0.59420) (1.18781,0.64706) (nan,nan)
(1.18781,0.64706) (1.17575,0.66250) (nan,nan)
(1.17575,0.66250) (1.16000,0.68047) (nan,nan)
(1.29042,-0.76471) (1.28000,-0.79269) (nan,nan)
(1.29042,-0.76471) (1.29563,-0.74938) (nan,nan)
(1.29563,-0.74938) (1.32817,-0.64706) (nan,nan)
(1.32817,-0.64706) (1.34680,-0.58157) (nan,nan)
(1.34680,-0.58157) (1.36035,-0.52941) (nan,nan)
(1.36035,-0.52941) (1.38313,-0.42831) (nan,nan)
(1.38313,-0.42831) (1.38640,-0.41176) (nan,nan)
(1.38640,-0.41176) (1.40000,-0.32777) (nan,nan)
(1.40000,0.16022) (1.39797,0.17448) (nan,nan)
(1.39797,0.17448) (1.39759,0.17647) (nan,nan)
(1.39759,0.17647) (1.37565,0.27024) (nan,nan)
(1.37565,0.27024) (1.36875,0.29412) (nan,nan)
(1.36875,0.29412) (1.34686,0.35967) (nan,nan)
(1.34686,0.35967) (1.32624,0.41176) (nan,nan)
(1.32624,0.41176) (1.31224,0.44337) (nan,nan)
(1.31224,0.44337) (1.28000,0.50622) (nan,nan)
(1.40543,-0.29412) (1.40000,-0.32777) (nan,nan)
(1.40543,-0.29412) (1.40620,-0.28804) (nan,nan)
(1.40620,-0.28804) (1.41666,-0.17647) (nan,nan)
(1.41666,-0.17647) (1.41765,-0.15916) (nan,nan)
(1.41765,-0.15916) (1.41978,-0.05882) (nan,nan)
(1.41978,-0.05882) (1.41946,-0.03974) (nan,nan)
(1.41946,-0.03974) (1.41383,0.05882) (nan,nan)
(1.41383,0.05882) (1.41263,0.07121) (nan,nan)
(1.41263,0.07121) (1.40000,0.16022) (nan,nan)
};
\addlegendentry{$z=3$}
\addplot[no marks,thin,orange] coordinates {
(1.40000,0.94376) (1.37159,0.97215) (nan,nan)
(1.37159,0.97215) (1.34038,1.00000) (nan,nan)
(1.45580,-1.00000) (1.50077,-0.90121) (nan,nan)
(1.50077,-0.90121) (1.50903,-0.88235) (nan,nan)
(1.50903,-0.88235) (1.52000,-0.85551) (nan,nan)
(1.52000,0.80434) (1.48634,0.84936) (nan,nan)
(1.48634,0.84936) (1.45851,0.88235) (nan,nan)
(1.45851,0.88235) (1.43112,0.91287) (nan,nan)
(1.43112,0.91287) (1.40000,0.94376) (nan,nan)
(1.55702,-0.76471) (1.52000,-0.85551) (nan,nan)
(1.55702,-0.76471) (1.57884,-0.70702) (nan,nan)
(1.57884,-0.70702) (1.60033,-0.64706) (nan,nan)
(1.60033,-0.64706) (1.63841,-0.53097) (nan,nan)
(1.63841,-0.53097) (1.63889,-0.52941) (nan,nan)
(1.63889,-0.52941) (1.64000,-0.52536) (nan,nan)
(1.64000,0.59184) (1.62351,0.63089) (nan,nan)
(1.62351,0.63089) (1.61558,0.64706) (nan,nan)
(1.61558,0.64706) (1.58265,0.70848) (nan,nan)
(1.58265,0.70848) (1.54819,0.76471) (nan,nan)
(1.54819,0.76471) (1.53711,0.78148) (nan,nan)
(1.53711,0.78148) (1.52000,0.80434) (nan,nan)
(1.67119,-0.41176) (1.64000,-0.52536) (nan,nan)
(1.67119,-0.41176) (1.68095,-0.37162) (nan,nan)
(1.68095,-0.37162) (1.69781,-0.29412) (nan,nan)
(1.69781,-0.29412) (1.71060,-0.22490) (nan,nan)
(1.71060,-0.22490) (1.71819,-0.17647) (nan,nan)
(1.71819,-0.17647) (1.72890,-0.08932) (nan,nan)
(1.72890,-0.08932) (1.73168,-0.05882) (nan,nan)
(1.73168,-0.05882) (1.73706,0.03633) (nan,nan)
(1.73706,0.03633) (1.73754,0.05882) (nan,nan)
(1.73754,0.05882) (1.73616,0.15310) (nan,nan)
(1.73616,0.15310) (1.73490,0.17647) (nan,nan)
(1.73490,0.17647) (1.72711,0.26187) (nan,nan)
(1.72711,0.26187) (1.72277,0.29412) (nan,nan)
(1.72277,0.29412) (1.71072,0.36345) (nan,nan)
(1.71072,0.36345) (1.69997,0.41176) (nan,nan)
(1.69997,0.41176) (1.68768,0.45851) (nan,nan)
(1.68768,0.45851) (1.66517,0.52941) (nan,nan)
(1.66517,0.52941) (1.65860,0.54765) (nan,nan)
(1.65860,0.54765) (1.64000,0.59184) (nan,nan)
};
\addlegendentry{$z=4$}
\addplot[no marks,thin,red] coordinates {
(1.69716,-1.00000) (1.74914,-0.89300) (nan,nan)
(1.74914,-0.89300) (1.75415,-0.88235) (nan,nan)
(1.75415,-0.88235) (1.76000,-0.86914) (nan,nan)
(1.80636,-0.76471) (1.76000,-0.86914) (nan,nan)
(1.80636,-0.76471) (1.83817,-0.68807) (nan,nan)
(1.83817,-0.68807) (1.85449,-0.64706) (nan,nan)
(1.85449,-0.64706) (1.88000,-0.57780) (nan,nan)
(1.88000,0.88424) (1.83529,0.95617) (nan,nan)
(1.83529,0.95617) (1.80467,1.00000) (nan,nan)
(1.89785,-0.52941) (1.88000,-0.57780) (nan,nan)
(1.89785,-0.52941) (1.90671,-0.50322) (nan,nan)
(1.90671,-0.50322) (1.93590,-0.41176) (nan,nan)
(1.93590,-0.41176) (1.95783,-0.33546) (nan,nan)
(1.95783,-0.33546) (1.96884,-0.29412) (nan,nan)
(1.96884,-0.29412) (1.99525,-0.18112) (nan,nan)
(1.99525,-0.18112) (1.99623,-0.17647) (nan,nan)
(1.99623,-0.17647) (2.00000,-0.15517) (nan,nan)
(2.00000,0.56984) (1.98443,0.63180) (nan,nan)
(1.98443,0.63180) (1.97987,0.64706) (nan,nan)
(1.97987,0.64706) (1.95482,0.72042) (nan,nan)
(1.95482,0.72042) (1.93735,0.76471) (nan,nan)
(1.93735,0.76471) (1.92010,0.80402) (nan,nan)
(1.92010,0.80402) (1.88110,0.88235) (nan,nan)
(1.88110,0.88235) (1.88072,0.88306) (nan,nan)
(1.88072,0.88306) (1.88000,0.88424) (nan,nan)
};
\addlegendentry{$z=5$}
\addplot[no marks,thin,brown] coordinates {
(1.92689,-1.00000) (1.97496,-0.90690) (nan,nan)
(1.97496,-0.90690) (1.98732,-0.88235) (nan,nan)
(1.98732,-0.88235) (2.00000,-0.85568) (nan,nan)
};
\addlegendentry{$z=6$}
\addplot[only marks,mark=*,black,forget plot] coordinates {(-0.378695,-0.464574)(2,-1)};
\end{axis}\end{tikzpicture}
\caption{Линии уровня $f_3$ и точки глобального минимума и максимума.
Изолинии построены линейной интерполяцией по треугольной сетке.}
\end{figure}

\subsection*{E.2. Точная область значений}
По $y$ выделим полный квадрат:
\begin{equation}\label{eq:complete}
f_3(x,y)=x^2+\left(y+\frac{\cos x}{2}\right)^2
+\sin x-\frac{\cos^2x}{4}.
\end{equation}
При любом $x\in[-1,2]$ точка $y=-\cos x/2$ лежит в $[-1,1]$.
Значит, для минимума на брусе остаётся минимизировать
\[
\varphi(x)=x^2+\sin x-\frac{\cos^2x}{4},\qquad
g(x)=\varphi'(x)=2x+\cos x+\frac12\cos x\sin x.
\]
Так как
\[
g'(x)=2-\sin x+\frac12\cos2x\geq\frac12>0,
\quad g(-1)<0,\quad g(2)>0,
\]
у $g$ ровно один корень. Бисекция даёт
\[
x_*\approx-0.378695393772,\qquad
y_*=-\frac{\cos x_*}{2}\approx-0.464573875151,
\]
\[
m=f_3(x_*,y_*)\approx-0.4421272966.
\]
Для проверки максимума заметим, что
\[
f_{xx}=2-\sin x-y\cos x\geq2-\sqrt2>0,
\quad f_{yy}=2,\quad f_{xy}=-\sin x.
\]
Определитель гессиана не меньше $3-2\sqrt2>0$.
Функция строго выпукла на брусе, поэтому максимум достигается
среди его вершин. В них значения равны
\[
\begin{array}{c|cccc}
(x,y)&(-1,-1)&(-1,1)&(2,-1)&(2,1)\\\hline
f_3(x,y)&2-\sin1-\cos1&2-\sin1+\cos1&
5+\sin2-\cos2&5+\sin2+\cos2
\end{array}
\]
Наибольшее из них $M=5+\sin2-\cos2\approx6.3254442634$.
Непрерывность на связном брусе даёт
\begin{equation}\label{eq:range3}
\boxed{\ran(f_3,X_3)=[m,M]\approx[-0.442127297,6.325444263]}.
\end{equation}
Бисекция используется только для вычисления единственного доказанно
существующего корня. Случайная оптимизация и предположение об
унимодальности неизвестной функции не требуются.
Дополнительно значения $f_3$ проверяются на контрольной сетке
из $101\times61$ точек. Эта проверка подтверждает согласованность
численного результата с найденными границами и дополняет
аналитическое доказательство глобальных экстремумов.

\subsection*{E.3. Естественное расширение исходной записи}
При принятом повторном умножении
\[
X_{3,1}^2=[-2,4],\quad X_{3,2}^2=[-1,1],\quad
\sin X_{3,1}=[-\sin1,1],\quad
\cos X_{3,1}=[\cos2,1].
\]
Поскольку $X_{3,2}\cos X_{3,1}=[-1,1]$, получаем
\begin{equation}\label{eq:nat3}
F_{3,E3}=[-4-\sin1,7]\approx[-4.841470985,7].
\end{equation}
Нижняя граница содержит, в частности, отрицательные произведения
двух независимых копий аргумента в квадратах.

\subsection*{E.4. Естественное расширение эквивалентной записи}
В выражении~\eqref{eq:complete} используем элементарное расширение
квадрата $Q$ вместо повторного умножения. Это оговорённая часть
вычислительной записи. Имеем
\[
Q(X_{3,1})=[0,4],\quad Q(\cos X_{3,1})=[0,1],
\quad Q\left(X_{3,2}+\tfrac12\cos X_{3,1}\right)=[0,9/4].
\]
Отсюда
\begin{equation}\label{eq:equiv3}
F_{3,E4}=[-\sin1-1/4,29/4]
\approx[-1.091470985,7.25].
\end{equation}
Для разделения двух эффектов можно заменить квадраты на $Q$
непосредственно в исходной записи: получится
$[-1-\sin1,7]\approx[-1.841470985,7]$ с расстоянием $1.399343688$.
Значит, выигрыш~\eqref{eq:equiv3} обусловлен как корректным учётом
одной и той же переменной внутри квадрата, так и выделением квадрата
по $y$. Верхняя граница после преобразования даже увеличилась,
но максимальная погрешность обеих границ уменьшилась.

\subsection*{E.5. Дифференциальная центрированная форма}
Частные производные:
\[
f_x=2x+\cos x-y\sin x,\qquad f_y=2y+\cos x.
\]
Для фиксированного $y$ первая производная $f_x$ возрастает по $x$,
поскольку $f_{xx}\geq2-\sqrt2>0$. По $y$ она аффинна.
Проверка концов даёт точный диапазон
\[
D_x=[-2+\cos1-\sin1,\ 4+\cos2+\sin2]
\approx[-2.301168679,4.493150590].
\]
Для второй производной слагаемые разделяются по переменным:
\[
D_y=[-2+\cos2,3]\approx[-2.416146837,3].
\]
Вычислим
\begin{equation}\label{eq:mv3}
F_{{\rm mv},3}(X_3,c)=f_3(c)+D_x(X_{3,1}-c_x)+D_y(X_{3,2}-c_y)
\end{equation}
для середины $c=(1/2,0)$ и точки $c=(0,0)$.

\begin{center}
{\small
\begin{tabular}{ccc}
\toprule
Центр & $f_3(c)$ & $F_{{\rm mv},3}$\\\midrule
$(1/2,0)$ & $0.729426$ & $[-9.0103,\,10.469151]$\\
$(0,0)$ & $0$ & $[-7.602337,\,11.986301]$\\
\bottomrule
\end{tabular}}
\end{center}

В отличие от одномерных примеров здесь по расстоянию лучше
оказался центр $(0,0)$, а не середина.

\subsection*{E.6. Бицентрированная среднезначная форма}
Правило~\eqref{eq:baumann} применяется отдельно к $D_x$ и $D_y$:
\[
p_x\approx0.3226197982,\qquad p_y\approx0.1077986216.
\]
Два центра на исходном брусе имеют координаты
\[
c^-\approx(0.016070303,-0.107798622),\qquad
c^+\approx(0.983929697,0.107798622).
\]
Соответствующие интервалы:
\[
F_{{\rm mv},3}(X_3,c^-)\approx[-7.321797309,12.157654461],
\]
\[
F_{{\rm mv},3}(X_3,c^+)\approx[-10.365378693,9.114073078].
\]
Пересечение даёт
\[
F_{3,E6}\approx[-7.321797309,9.114073078].
\]
Оно улучшает обе рассмотренные одиночные среднезначные оценки
по расстоянию до точного диапазона, но уступает естественному
расширению~\eqref{eq:nat3}.

\subsection*{E.7. Дробление бруса}
Середины координатных интервалов равны $1/2$ и $0$.
При разбиении по $x$ используем $[-1,1/2]\times[-1,1]$ и
$[1/2,2]\times[-1,1]$, по $y$ --- $[-1,2]\times[-1,0]$ и
$[-1,2]\times[0,1]$. На каждом подбрусе вычисляется именно
исходная естественная форма E.3 с повторным умножением.
Оболочка двух результатов определяется как
\[
\hull([a,b]\cup[c,d])=[\min(a,c),\max(b,d)].
\]

\begin{center}
{\small
\begin{tabular}{lccc}
\toprule
Разбиение & Первый подбрус & Второй подбрус & Оболочка\\\midrule
По $x$ & $[-3.341471,\,3.479426]$ & $[-1.148157,\,6.877583]$ & $[-3.341471,\,6.877583]$\\
По $y$ & $[-3.841471,\,6.416147]$ & $[-3.257618,\,7]$ & $[-3.841471,\,7]$\\
\bottomrule
\end{tabular}}
\end{center}

Оба разбиения сужают исходную оценку. Деление по $x$ улучшает
сразу обе границы и даёт меньшую погрешность. При делении по $y$
улучшается только нижняя граница общей оболочки.
Разница объясняется структурой вычисляемого выражения: аргумент $x$
участвует в квадрате, синусе и косинусе, поэтому сужение его
интервала уточняет сразу несколько зависимых слагаемых.
Разбиение по $y$ не изменяет диапазоны тригонометрических функций.
Полученный выигрыш относится к данному брусу и записи функции;
он не задаёт универсального правила выбора координаты дробления.

\subsection*{E.8. Покоординатная оценка Липшица}
Точные диапазоны производных позволяют взять
\[
L_1=\max|D_x|=4+\cos2+\sin2\approx4.493150590,
\qquad L_2=\max|D_y|=3.
\]
Проверяем неравенство из задания для исходной формы:
\[
\rad F_{3,E3}=\frac{11+\sin1}{2}\approx5.920735492,
\]
\[
L_1\rad X_{3,1}+L_2\rad X_{3,2}
=\frac32(4+\cos2+\sin2)+3\approx9.739725885.
\]
Таким образом, $5.920735492\leq9.739725885$.
Для данного бруса проверяемое неравенство выполнено, однако диапазоны
производных точечной функции сами по себе не дают универсальной
оценки полуширины произвольной интервальной записи.
Дополнительно по дереву исходного выражения можно взять
\[
L_{{\rm expr},x}=2|X_{3,1}|+\max|\cos X_{3,1}|
+|X_{3,2}|\max|\sin X_{3,1}|=4+1+1=6,
\]
\[
L_{{\rm expr},y}=2|X_{3,2}|+|\cos X_{3,1}|=3.
\]
Они дают обоснованную для этой записи границу $6\cdot1.5+3=12$,
которая также превышает фактическую полуширину.

\subsection*{E.9. Точность оценок и выводы}
В таблице сопоставляются оценки на всём исходном брусе.
Для разбиений используются оболочки, поскольку каждый отдельный
подбрус охватывает лишь часть исходной области.

\begin{center}
{\small
\begin{tabular}{lcc}
\toprule
Метод & Внешняя оценка & $\dist(F,\ran f_3)$\\\midrule
E.3, исходная запись & $[-4.841471,\,7]$ & $4.399344$\\
E.4, полные квадраты & $[-1.091471,\,7.25]$ & $0.924556$\\
E.5, центр $(1/2,0)$ & $[-9.0103,\,10.469151]$ & $8.568173$\\
E.5, центр $(0,0)$ & $[-7.602337,\,11.986301]$ & $7.16021$\\
E.6, два центра & $[-7.321797,\,9.114073]$ & $6.87967$\\
E.7, дробление по $x$ & $[-3.341471,\,6.877583]$ & $2.899344$\\
E.7, дробление по $y$ & $[-3.841471,\,7]$ & $3.399344$\\
\bottomrule
\end{tabular}}
\end{center}

Наименьшее расстояние среди этих оценок даёт выражение
с выделенными квадратами E.4: примерно $0.924556$.
При этом оно не вложено в E.3, так как имеет большую верхнюю границу.
Улучшение по расстоянию не означает вложенность интервалов.
Из двух направлений дробления эффективнее оказалось $x$:
расстояние уменьшилось с $4.399344$ до $2.899344$,
тогда как при дроблении по $y$ оно равно $3.399344$.

Среднезначная форма на широком брусе суммирует независимые
диапазоны частных производных и отклонений от центра, что создаёт
дополнительную переоценку. Она является строгим включением при
точных интервалах производных, а не приближением первым членом
ряда с отброшенным остатком. Причина её ширины --- утрата связи
между множителями, а не отсутствие учёта нелинейных членов.

\section*{Выводы}
\begin{enumerate}
\item Для $f_1$ и $f_2$ найдены точные области значений и вычислены
внешние интервальные оценки всеми методами B.1--B.5.
Для $f_3$ минимум установлен сведением к одной переменной,
а максимум --- проверкой вершин выпуклой функции на брусе.

\item Алгебраическая запись влияет на переоценку. Для $f_1$
схема Горнера сужает результат с $[-25,29]$ до $[-25,2]$.
Для $f_2=e^x-2x-1$ группировка линейных слагаемых оставляет
оценку $[-4,e^2-1]$ без изменения. Поэтому преобразование
выражения не обязательно улучшает результат.

\item Дифференциальные центрированные формы чувствительны
к выбору центра. Среди проверенных центров середина даёт
наименьшее расстояние для обеих одномерных функций.
Для $f_3$ центр $(0,0)$ оказался лучше середины $(1/2,0)$.
В наклонных формах для $f_1$ и $f_2$ лучший результат среди
рассмотренных центров получен в левом конце интервала.

\item Для $f_1$ наименьшее расстояние среди всех рассмотренных
одномерных оценок даёт наклонная форма при $c=0$: $2.75$.
Для $f_2$ лучшей оказалась бицентрированная форма с расстоянием
около $1.559171$. Для $f_3$ пересечение форм Бауманна улучшает
проверенные одиночные среднезначные формы, но уступает
естественному расширению исходной записи.

\item Константу Липшица точечной функции следует отличать
от константы её вычислительного выражения. Для исходной записи
$f_1$ значение $L_{f_1}=9$ не ограничивает полуширину
естественной оценки. Константа $L_{1,B1}=45$, найденная
по дереву Канторовича, даёт выполненное неравенство $27\leq67.5$.

\item Для естественной оценки $f_2$ полуширина составляет
примерно $5.194528$, а граница по дереву выражения --- $9.389056$.
Для $f_3$ проверена покоординатная оценка
$5.920735\leq9.739726$. Дополнительная граница по дереву
исходного выражения равна $12$ и также выполняется.

\item Дробление исходного бруса улучшает естественную оценку
$f_3$ в обоих направлениях. Более эффективным оказалось деление
по $x$: расстояние уменьшается с $4.399344$ до $2.899344$,
а по $y$ --- до $3.399344$. Наименьшее расстояние в двумерном
сравнении даёт запись с выделенными элементарными квадратами:
примерно $0.924556$.

\item Все рассчитанные внешние оценки содержат найденные
области значений с учётом указанной точности численных вычислений.
Программные проверки включают охват диапазонов, вложение
наклонных форм в соответствующие среднезначные оценки
и выполнение рассмотренных неравенств для полуширин.
\end{enumerate}

\section*{Источники}
\begin{enumerate}
\item \textit{Лабораторная работа №2 по интервальному анализу.}
Задание от 17 сентября 2026 г., вариант 23.\\
Файл: \texttt{task\_2\_new.pdf}.

\item Баженов А.~Н., Карпова А.~А.
\textit{Интервальный анализ в примерах.}\\
Санкт-Петербург, 2024. Разделы 1.3.1--1.3.4.
\end{enumerate}

\noindent\textbf{Репозиторий с исходным кодом:}\\[2mm]
\url{https://github.com/Jujun1or/intan/tree/lab2}
\end{document}
